\section{Method}
\label{sec:method}

Figure~\ref{fig:pipeline} gives an overview. The benchmark and the
Azerbaijan prototype use the same three-channel SAR representation.
The benchmark scores the raw output of DeepLabV3+. The Azerbaijan path
runs a different network on scene tiles and applies two downstream
exclusions to the resulting water mask.

\begin{figure}[!htbp]
\centering
\begin{tikzpicture}[
  node distance=3mm and 4mm,
  stage/.style={draw=black!55, fill=black!4, rounded corners=1pt,
    align=center, font=\small, minimum height=10mm, text width=21mm},
  every edge/.style={draw, -{Stealth[length=2mm]}},
]
\node[stage] (study) {DeepLabV3+\\ResNet-50};
\node[stage,right=of study] (test) {Raw water\\test masks};
\node[stage,right=of test] (iou) {Pooled\\water IoU};
\node[anchor=south west,inner sep=0pt,font=\small\bfseries] at
  ([yshift=1.5mm]study.north west) {Sen1Floods11 benchmark};
\path (study) edge (test) (test) edge (iou);
\node[stage,below=10mm of study] (model) {SegFormer-B4\\scene tiles};
\node[stage,right=of model] (raw) {Raw stitched\\water mask};
\node[stage,right=of raw] (filter) {JRC then\\DEM filter};
\node[stage,right=of filter] (map) {Azerbaijan\\map product};
\node[anchor=south west,inner sep=0pt,font=\small\bfseries] at
  ([yshift=1.5mm]model.north west) {Azerbaijan mapping prototype};
\path (model) edge (raw) (raw) edge (filter) (filter) edge (map);
\end{tikzpicture}

\caption{Benchmark and prototype paths. Both use VV, VH and VV$-$VH.
Water IoU for the event-held-out comparison is computed on raw
DeepLabV3+ output. The Azerbaijan prototype filters the raw stitched
SegFormer-B4 water mask with JRC occurrence and DEM-based terrain and
coverage information before rendering the map.}
\label{fig:pipeline}
\end{figure}

\subsection{Input Representation}
\label{sec:input}

The network input is a three-channel SAR image of size $512\times512$
with the channels VV (vertical--vertical), VH (vertical--horizontal),
and their ratio, computed as the difference VV$-$VH in dB. The
polarizations are complementary. VV carries the surface backscatter,
and VH helps to separate land from water, because smooth water
backscatters much less than surrounding land. The ratio channel emphasizes
the difference between the two polarizations. Using three channels also
matches the input format expected by ImageNet-pretrained encoders
\cite{deng2009}, so the pretrained first layer can be used without
modification.

Each channel is clipped to a fixed dB range and then linearly rescaled to
$[0,1]$: VV to $[-25,0]$ dB, VH to $[-35,0]$ dB, and VV$-$VH to
$[0,18]$ dB. The exploratory analysis used to set these ranges drew from
the full hand catalog. Its reconstructed 80-chip sample includes 16
official test SAR inputs, three Bolivia holdout inputs and all 11
hand-labeled groups. The fixed ranges were shared by both comparison
arms and incorporate benchmark test covariates.

Wind roughens the water surface and raises its backscatter, so water can
resemble land. Conversely, recurrent water and radar shadow can both
look dark. This motivates the two downstream exclusions in
Sec.~\ref{sec:fusion}.

\subsection{Architecture}

We use DeepLabV3+ \cite{chen2018deeplabv3plus} with an ImageNet-initialized
ResNet-50 encoder \cite{he2016resnet} for the event comparison.
DeepLabV3+ is an encoder--decoder architecture for semantic segmentation.
Dilated convolutions and an Atrous Spatial Pyramid Pooling (ASPP) module
capture context at multiple scales, and a lightweight decoder refines
the result to sharpen object boundaries.

\subsection{Two-Stage Training}
\label{sec:training}

Hand-labeled data are scarce, so we train in two stages: pretraining
on chips whose labels were derived automatically from Sentinel-2 imagery,
followed by fine-tuning on manually annotated chips. The first stage
exposes the model to a large but noisy dataset, and the second adapts it
to expert labels. The named held-out event is excluded from both stages
in the reconstructed LOCO manifests (Sec.~\ref{sec:protocol}). We compare
with a single-stage baseline trained on hand labels in
Sec.~\ref{sec:ablations}.

\paragraph{Optimization.}
Both stages use Adam with a constant learning rate and no weight decay
or scheduler, and Dice loss except where noted. The effective batch
size is 8 (batch size 2 with 4 gradient-accumulation steps), and training
uses mixed precision (fp16). Pretraining uses a learning rate of
$3\times10^{-4}$ for up to 20 epochs with early stopping (patience 5).
Fine-tuning uses $10^{-4}$ for up to 50 epochs, keeping the checkpoint
with the best validation IoU. Training augmentation is horizontal flip,
vertical flip and right-angle rotation, each applied independently
with probability 0.5. Validation and test chips are not augmented.

\subsection{Loss Function}

With water in roughly one in ten labeled training pixels, the loss must
cope with strong class imbalance. We compared 11 variants, including
binary cross-entropy (BCE) with positive weights of 1, 4, 8.4 and 13,
and overlap-based objectives. We adopt the Dice loss \cite{milletari2016},
shown without its numerical stabilizer:
\begin{equation}
\mathcal{L}_{\mathrm{Dice}}
=1-\frac{2\sum_i p_i g_i}{\sum_i p_i+\sum_i g_i},
\label{eq:dice}
\end{equation}
where $p_i\in[0,1]$ is the predicted water probability and
$g_i\in\{0,1\}$ the label on valid pixel $i$. Dice was within 0.001
validation IoU of the leading Tversky variant in the short ResNet-34
sweep, and needs no additional class-imbalance hyperparameter
(Sec.~\ref{sec:ablations}).

\subsection{Post-Processing Fusion}
\label{sec:fusion}

\paragraph{Recurrent water (JRC).}
We use JRC Global Surface Water occurrence \cite{pekel2016high} and
subtract predicted positives where valid historical occurrence is at
least 20\%. This threshold is a recurrence proxy. It can remove genuine
event-date water as well as persistent water. Missing JRC pixels trigger
no JRC exclusion.

\paragraph{Layover and shadow (DEM).}
Radar terrain distortions depend on the viewing geometry and local
slope \cite{vollrath2020}. On a north-up raster, columns increase
eastward and rows increase southward. Let $s$ be the terrain slope
magnitude and $\alpha$ the downhill aspect, clockwise from north.
The sensor-to-ground azimuth $\varphi$ is nominally $90^\circ$ for
ascending passes and $270^\circ$ for descending passes. For nominal
incidence $\theta_i$, the signed angle in the range plane is
\begin{equation}
\theta_{\mathrm{loc}}=\theta_i+
\arctan\!\bigl(\tan s\cos(\varphi-\alpha)\bigr).
\label{eq:local-incidence}
\end{equation}
Terrain rising along the look direction faces the sensor. Its
downhill projection along that direction is negative, reducing
local incidence. Pixels with $\theta_{\mathrm{loc}}\leq0^\circ$ are
flagged for layover and those with $\theta_{\mathrm{loc}}\geq90^\circ$
for shadow. Both scene replays use SRTM1 \cite{farr2007}, a nominal
incidence of $37.55^\circ$ (the midpoint of $29.1^\circ$ and $46.0^\circ$),
and the eastward look approximation. Missing DEM is excluded along
with a one-pixel dilation of the coverage gaps.

\paragraph{Order of operations and evaluation.}
The JRC mask is subtracted from predicted water first, and the
terrain/coverage mask is applied second. Sen1Floods11 labels include
persistent water present on the event date, so the benchmark scores raw
water masks before these exclusions. We examine the filtered map as a
separate downstream product in Sec.~\ref{sec:deployment}.
